\cotefile{57}
\scanpage{89}
Donnons-nous une section \deleted{\unsure{de}} $\unsure{\zeta|\xi}$ de $\underline{\mathrm{Pic}}_{X/S}^{\omega}$ sur $S$, d'où par image inverse un élément de $H^{1}(S,\underline{\omega}_{X/S}/\operatorname{Im}f_{*}\mathcal{O}_{X}^{*}))$, que nous \unsure{voudrions|voulons} calculer. Pour simplifier, nous allons supposer que $\zeta$ provient d'un faisceau inversible $\underline{L}$ sur $X$. Alors on trouve l'image \deleted{\unsure{de|est}} l'élément de $H^{1}(S,\underline{\omega}_{X/S})\overset{\mathrm{Ker}}{\simeq}(H^{1}(X,\Omega^{1}_{X/S})\unsure{\to|-}H^{0}(S,R^{1}f_{*}\unsure{\Omega|d\Omega|\underline{\Omega}}^{1}_{X/S}))$ défini par $\underline{L}$.

\unsure{1}

% archive sheet number: 88

\scanpage{90}
\mspage{\handwritten{57}}
limit ourselves to the case $g=2$. We are interested in $M(\Gamma_{2}(\unsure{\lambda|l}))$ for $\lambda=1,2$, because we know $S(\Gamma_{2}(\lambda))$ explicitly for the same values of $\lambda$ (8). We keep in mind that they are defined as quotient varieties and not as monoidal transforms. We take a point $t_{o}$ of $\underline{S}_{1}$. Since $t_{o}$ is not a fixed point of any element of $\Gamma_{1}(2)$ different from $\pm1_{2}$, the analytic local ring $\unsure{O|0}$ of $M(\Gamma_{2}(2))$ at the image point of $t_{o}$ consists of convergent series of the form
\[
\sum_{k=\unsure{o|0}}^{\infty}\theta_{k}(t,z)e(kw/2)\,.
\]
The coefficient $\theta_{k}(t,z)$ is holomorphic in $V\handwritten{\times}\mathbb{C}$ for some $V$ and $z\longrightarrow\theta_{k}(t,2z)$ is a theta-function of order $\handwritten{4}k$, which is even in the sense it is invariant under $z\longrightarrow-z$. This follows from Theorem 1. Moreover, the ideal $I$ consists of convergent series which start from $k=1$. Furthermore, there exists a base $(\theta_{\unsure{l|1}}(t_{o},z))$ over $\mathbb{C}$ of the vector space of even theta-functions of order four such that each $\theta_{l}(t,z)$ is holomorphic in some $V\handwritten{\times}\mathbb{C}$ and such that the correspondence $z\longrightarrow(\theta_{l}(t,z))$ gives rise to a complex-analytic embedding of the quotient variety of $T_{1}(t)$ by its involution $z\longrightarrow-z$, which is $P_{1}(\mathbb{C})$, as a plane curve of order two. Therefore, if we take the monoidal transform of $S(\Gamma_{2}(2))$ along its boundary in the neighborhood of the image point of $t_{o}$, we get the quotient variety of the monoidal transform of $S(\Gamma_{2}(\lambda))$, for any even $\lambda$ at least equal to 3, along its boundary in the neighborhood of the image point of $t_{o}$. In the same way, we can conclude that the morphism $M(\Gamma_{2}(1))\longrightarrow S(\Gamma_{2}(1))$ is a monoidal transformation in the neighborhood of the image point of $t_{o}$ provided $t_{o}$ is not a fixed point of any element of $\Gamma_{1}(1)$ different

% archive sheet number: 89

\scanpage{91}
\uline{NB} Si on a un hom
\[
u:B\longrightarrow\underline{\mathrm{Pic}}_{X/S}
\]
avec $B$ un schéma abélien, et si $\underline{\Omega}^{1}_{X/S}$ coh. plat $/S$ en dim $0$, alors $u$ se factorise en $B\longrightarrow\underline{\mathrm{Pic}}_{X/S}^{\omega}$. En particulier, si $\underline{\mathrm{Pic}}_{X/S}^{00}$ existe, alors on a
\[
\underline{\mathrm{Pic}}_{X/S}^{00}\subset\underline{\mathrm{Pic}}_{X/S}^{\omega}.
\]

Notons que la \uline{donnée} de $u$ équivaut : \margin{sous réserve d'existence de $B^{*}$ et de possibilité de descente (p.ex. si $B/S$ est projectif)} la donnée d'un torseur $A$ sous \deleted{$A$} le schéma dual $B^{*}$ de $B$, et d'un morphisme
\[
X\longrightarrow A
\]
\[
[\,\text{d'où}\qquad B\simeq\underline{\mathrm{Pic}}^{0}_{A/S}\subset\underline{\mathrm{Pic}}_{A/S}\longrightarrow\underline{\mathrm{Pic}}_{X/S}\quad\ldots\,]
\]
On trouve alors, par fonctorialité\margin{fonctorialité doit être explicitée avec un peu de soin !} un diagramme
\diagram{\begin{tikzcd}[ampersand replacement=\&,column sep=small,row sep=small]
0\arrow[r,dash] \& \underline{\omega}^{\unsure{*}}_{X/S}\arrow[r] \& \mathcal M_{X/S}\arrow[r] \& \underline{\mathrm{Pic}}^{\omega}_{X/S}\arrow[r] \& 0\\
0\arrow[r] \& \underline{\omega}_{A/S}\arrow[u]\arrow[r] \& \mathcal M_{A/S}\arrow[u]\arrow[r] \& \underline{\mathrm{Pic}}^{\omega}_{A/S}\arrow[u]\arrow[r,dash] \& 0\\
0\arrow[r,dash] \& \underline{\omega}_{A/S}\arrow[u,equal]\arrow[r] \& \deleted{\unsure{M_{B/S}|\underline E}}\mathcal M^{0}_{A/S}\arrow[u]\arrow[r] \& B\arrow[u,hook']\arrow[r] \& 0\\
\& V(R^{\unsure{1}}f_{B*}(\underline{\mathcal{O}}_{B}))\arrow[u,"\simeq" description,phantom] \& \& \underline{\mathrm{Pic}}^{0}_{A/S}\arrow[u,equal]
\end{tikzcd}}

Ici $\deleted{\unsure{M_{B/S}|\underline E_{B^*/S}}}\mathcal M^{0}_{A/S}$ est défini comme image inverse de $\mathcal M_{A/S}$ \inserted{par $B\hookrightarrow\underline{\mathrm{Pic}}^{\omega}_{A/S}$} [ou aussi de $\mathcal M_{X/S}$ par $B\longrightarrow\underline{\mathrm{Pic}}^{\omega}_{X/S}$ \inserted{si $\underline{\omega}_{A/S}\simeq\underline{\omega}_{X/S}$}]. Je vou-

% archive sheet number: 90

\scanpage{92}
drais l'interpréter comme l'extension \uline{universelle} (de Serre) de $B$ par un fibré vectoriel.

Notons d'abord que le schéma $B^{*}$ des translations de $A$ opère \unsure{triv\textsuperscript{t}|trivialt|trivialement} sur l'extension $\deleted{\unsure{M_{A/S}|\underline E_{A/S}}}\inserted{\mathcal M_{A/S}^{0}\deleted{\unsure{=|\underline E}}}$ (car il\deleted{s} opère \unsure{triv\textsuperscript{t}|trivialement} sur $\underline{\mathrm{Pic}}^{0}_{A/S}$ et sur $\underline{\omega}_{A/S}$\deleted{)}, et $\unsure{\forall|\mathrm{tt}}$ hom de $B$ dans $\underline{\omega}_{X/S}$ est trivial). Donc par théorie de descente on voit que $\mathcal M^{0}_{A/S}$ est can isom à $\mathcal M^{0}_{B^{*}/S}$, donc c'est un invariant qui ne dépend que du \deleted{\unsure{le}} schéma abélien $B$ sur $S$.

Pour calculer la classe de l'extension obtenue de $B$ par $\underline{\omega}_{A/S}$ comme élément de $H^{0}(S,\underbrace{R^{1}f_{B*}^{\deleted{\unsure{1}}}(\underline{\omega}_{A/S}\otimes_{\mathcal O_S}\mathcal O_B)}_{\unsure{(R^{1}f_{B*}\mathcal O_{B})|R^{1}f_{B*}(\mathcal O_{B})|R^{1}f_{B*}(\underline{\mathcal O}_{B})}^{\vee}\otimes R^{1}f_{B*}(\mathcal O_{B})})$, on regarde \inserted{par la recette générale} le faisceau inversible \unsure{canonique|unique|de Poincaré} sur $P=A\times_{S}B$, et on écrit
\diagram{\begin{tikzcd}[ampersand replacement=\&,column sep=small,row sep=small]
A \& \& B \& A\arrow[d,equal] \&[-0.8em] \times_{S}B\\
\& S\arrow[ul,"f_A"']\arrow[ur,"f_B",dash] \& \& B^{*} \& {}
\end{tikzcd}}
l'obstruction à mettre dessus une connexion relativement à $B$, i.e. la classe de cohomologie qui lui correspond dans $H^{1}(P,\underline{\Omega}^{1}_{P/B})$ [puis on localise sur $S$].

Calculons donc la classe de \unsure{cohomologie} de $\underline L$ dans\deleted{\unsure{à}} $H^{1}(A\times_{S}B,\underline{\Omega}^{1}_{A\times_{S}B/S})$\deleted{$\unsure{\in H|\text{de }L}$}, localisé en $S$.

% archive sheet number: 91

\scanpage{93}
\deleted{Cher Monsieur \unsure{Che}\illegible{1}}\\

% archive sheet number: 92

\scanpage{94}
\[
\begin{aligned}
\underline{H}^{11}(A\times_{S}B)&=\underline{H}^{11}(A/S)\oplus\underline{H}^{11}(B/S)\\
&\quad+\underline{H}^{1,0}(A/S)\otimes\underline{H}^{0,1}(B/S)+\underline{H}^{0,1}(A/S)\otimes\underline{H}^{1,0}(B/S)
\end{aligned}
\]
La classe $\unsure{\zeta|\xi}$ qu'on veut calculer a, pour des raisons de fonctorialité (regarder les inclusions $A\longrightarrow A\times_{S}B$, $B\longrightarrow A\times_{S}B$) des composantes nulles sur les deux premiers facteurs. Il reste à calculer sa classe dans le deuxième, on \unsure{trouve|devine|peut} ce qu'on \unsure{pense|fera|se faire}....

\begin{center}\rule{3cm}{0.35pt}\\[-0.6ex]\rule{4cm}{0.45pt}\end{center}

\deleted{Connexion \inserted{(absolue)} canonique sur $\underline{\mathcal M}_{X/S}$}

\deleted{\diagram{\begin{tikzcd}[ampersand replacement=\&,column sep=small]
X_{0}\arrow[r,dash]\arrow[d,dash] \& X'_{\unsure{1}}\arrow[d,dash] \& X''_{2}\arrow[dl,dash]\\
S_{0}\arrow[r,phantom,"\subset" description] \& S^{\unsure{\prime|' '}}
\end{tikzcd}}}
\margin{\deleted{dans le voisinage inf. du premier ordre de $S$ avec $\exists$ rétraction de $S$ sur $S_{0}$}}

\deleted{$\underline{W}_{X'_{1}/S'}\simeq\underline{W}_{X'_{2}/S'}$ canoniquement ??}

\deleted{\uline{Question équivalente} : Si on a un $L_{0}$ inversible sur $X_{0}$, muni d'une connexion\unsure{,} rel. à $S_{0}$, alors les prolongements (avec connexion) à $X'$ et à $X''$ se correspondent \unsure{biuniv.quement|biunivoquement}.}

$W_{\unsure{X/S}}$

$S_{0}\leftarrow S$

% archive sheet number: 93

\scanpage{95}
\mspage{\handwritten{56}}
an integer matrix. After this remark, we consider all sequences in $\underline{S}_{g}$ of the form $(t,z)\longrightarrow(t_{o},z_{o})$ and the normal coordinates of $\operatorname{Im}(w)$ tending to $-\infty$ satisfying the condition $y_{12}\geqq y_{13},\,y_{23}$. Then we certainly have $\xi\longrightarrow0$. Moreover, if we take $\theta_{\sigma_{12}}(t,z)$ such that we have $\theta_{\sigma_{12}}(t_{o},z_{o})\neq0$, the quotient $(H_{\sigma}/H_{\sigma_{12}})(t,z,w)$ is finite along the sequences for every integral, positive, non-degenerate matrix $\sigma$. Among these quotients, we find
\[
(H_{\sigma_{i3}}/H_{\sigma_{12}})(t,z,w)=(\theta_{\sigma_{i3}}/\theta_{\sigma_{12}})(t,z)(\xi_{i3}/\xi_{12})(1+\ldots)
\]
for $i=1,2$. Moreover $(\xi_{13}/\xi_{12},\xi_{23}/\xi_{12})$ can approach to any point of the product of two copies of a closed unit disc in $\mathbb C$ provided we take the limit $\operatorname{Im}(w)\longrightarrow\infty$ suitably. This implies that the fiber at the image point of $t_{o}$ of the monoidal transform of $S(\Gamma)$ along its boundary is of dimension one more than the fiber of $M(\Gamma)\longrightarrow S(\Gamma)$ at the image point of $t_{o}$. Therefore $M(\Gamma)$ is not the monoidal transform of $S(\Gamma)$ along its boundary. This completes the proof.

We note that the blowing up $M(\Gamma)\longrightarrow S(\Gamma)$ can be analyzed along the quasi projective varieties of the form $\Gamma_{o}\backslash\underline S_{g_{o}}$ for $g_{o}=g-1$, contained in the boundary. In fact, it is the partial desingularization of $\unsure{\underline S|S}(\Gamma)$ along the image of $\Gamma_{o}\backslash\underline S_{g_{o}}$, provided $\Gamma$ operates without fixed points on $\underline S_{g}$ (cf. 10). We also note that most of the groups in the commensurability class of $\operatorname{Sp}(g,\mathbb Z)$ has ``undesirable'' properties stated for $\Gamma_{g}(\unsure{\lambda|l},2\lambda)$. We shall not try to give a precise form to this statement.

We shall examine $\unsure{M|\underline M}(\Gamma_{g}(\lambda))$ for $\lambda=1,2$. Although a part of the following investigations can be carried out also for $g=3$, we shall

% archive sheet number: 94

\scanpage{96}
``Critères d'équivalence''

% archive sheet number: 95

\scanpage{97}
\begin{center}\fbox{Critères d'équivalence}\end{center}

$X$ variété projective \deleted{\unsure{non sing|irréductible} \unsure{normale}}\inserted{non sing} de dim $n$

$Y$ section hyperplane \deleted{\unsure{normale}} non sing

$X_{m}$ $m$.\textsuperscript{ème} voisinage infinitésimal de $X$

\diagram{\(\left\{\begin{array}{@{}l@{\quad}p{.60\linewidth}@{}}
\mathrm{Pic}(X)\to\mathrm{Pic}(Y) &
\(\left\{\begin{array}{@{}p{.55\linewidth}@{}}
par \uline{Lef}-Grauert $\downarrow$\\
inj si $n\geq3$, $H^{1}(Y,\mathcal O_Y(\unsure{m|-m}))=0$ pour $m>0$, surjectif si $n\geq4$, $H^{2}(Y,\mathcal O_Y(\unsure{-m|+m|m}))=0$ pour $m>0$\\
noyau un $p$-groupe fini si $n\geq3$ (nul si $Y$ \uline{général}), \fbox{noyau nul si $n\geq4$}\\
$\leftarrow$ Pic des schémas non \unsure{réduits}, \unsure{rel. dim|rel. à un|$X_{m}$}
\end{array}\right.\)\\
\mathrm{Pic}^{\tau}(X)\to\mathrm{Pic}^{\tau}(Y) &
inj. exc. $p$ si $n\geq2$, bij. exc $p$ si $n\geq3$, injectif si $n\geq2$ et $Y$ \uline{général}.\\
& $\nwarrow$ via $\pi_{1}$, Lef. Grauert\\
& $\swarrow$ via $\pi_{1}$, Lef. Grauert\\
\mathrm{Pic}^{0}(X)\to\mathrm{Pic}^{0}(Y) &
inj exc. $p$ si $n\geq2$, surjectif si $n\geq3$, bijectif si $n\geq2$ et $Y$ \uline{général}
\end{array}\right.\)}

\diagram{\(\begin{array}{@{}l@{\quad}p{.61\linewidth}@{}}
TN(X)\to TN(Y)& bij exc. $p$ si $n\geq3$ [Si $n\geq3$, le noyau est l'image \deleted{de} celui de $\mathrm{Pic}^{\tau}(X)\to\mathrm{Pic}^{\tau}(Y)$] \uline{injectif} si $Y$ \uline{général} et $n\geq3$, ou si $n\geq4$\\
LN(X)\to LN(Y)& injectif si $n\geq3$\\
N(X)\to N(Y)& noyau un $p$-groupe fini si $n\geq3$ [contenu exc. $p$ dans l'image du noyau \deleted{de} $\mathrm{Pic}(X)\to\mathrm{Pic}(Y)$, qui est un $p$-groupe fini.] \uline{nul} si on a $\unsure{\overline m}$ $n\geq4$ \uline{ou} $Y$ \uline{général}\\
\underset{m\ \text{grand}}{\mathrm{Pic}(X)\longrightarrow\mathrm{Pic}(X_{m})}& inj si $n\geq3$, bijectif si $n\geq4$
\end{array}\)}

\uline{Question} Surjectivité \inserted{(donc bijectivité)} de $\mathrm{Pic}(X)\to\mathrm{Pic}(Y)$ \deleted{\unsure{sans|pour}} pour $n\geq4$ \uline{sans conditions cohomologiques}, [\uline{équivaut} à celle de $\deleted{\unsure{T}}N(X)\to\deleted{\unsure{T}}N(Y)$, ou encore celles de $TN(X)\to TN(Y)$ et $LN(X)\to LN(Y)$]. Équivaut à celle de $\mathrm{Pic}(X_{m})\to\mathrm{Pic}(Y)$ pour $m$ grand. N.B. On a surjectivité mod $p^{m}$ (car. $p>0$)\inserted{(regarder les obstructions successives, sont annulées par $p$\ldots.)} \uline{donc pour} $n\geq4$, $\mathrm{Pic}(X)\to\mathrm{Pic}(Y)$ \uline{est surjectif exc.} $p$.
% slanted
\margin{Corollaire : $\rho(X)=\rho(Y)$ pour $n\geq4$ [$\rho(X)\leq\rho(Y)$ si $n\geq3$] \deleted{$\rho(X)\leq\rho(Y)+1$ si $n=2$ et $\unsure{Y}$ \uuline{général}]}}
\deleted{\unsure{Qua|Que}} En car. $0$, \uline{on peut appliquer la théorie Kählérienne}, qui prouve que \deleted{\illegible{1}} $N\deleted{\unsure{S}}(X)\overset{\sim}{\longrightarrow}N(Y)$, d'où $\mathrm{Pic}(X)\longrightarrow\mathrm{Pic}(Y)$.

Démonstration algébrique ???

% archive sheet number: 96

\scanpage{99}
\mspage{1}
\uuline{Théorème}\qquad $X/k$, $k$ corps \uuline{alg. clos}

$X$ \deleted{\uuline{intègre}}\inserted{\uline{normal} \uline{\unsure{irred|irré d|intègre}}} \uuline{projectif}, $X\subset\mathbb P^{r}_{k}$

\deleted{$\operatorname{prof}X\geq2$ (d'} $\dim X\geq2$ (d'où $\operatorname{prof}X\geq2$)

$G$ groupe algébrique \uuline{affine} \uuline{sans \unsure{composants|composantes}} $G_{m}$. (p.ex. un groupe alg. fini).

$K=\deleted{\unsure{\overline{\phantom{k(t_1,t_r)}}}}k(t_{1},\ldots,t_{r})$, $\overline K$ sa clôture algébrique

$Y_{\overline K}$ la section hyperplane générique, définie sur $\overline K$ :

\[
\text{Alors}\qquad H^{1}(X,G)\longrightarrow H^{1}(Y_{\overline K},G)\qquad\text{\uuline{est injectif}}
\]

La structure des groupes algébriques comm. nous ramène, aux cas simples $G_{a}$, \deleted{\unsure{$G_{m}$}} groupe fini séparable, \inserted{$\alpha_{p}$}, $\mu_{p}$. Donc on est ramené à prouver

% slanted
\margin{N.B. Généralisation au cas d'un morphisme donné $X\longrightarrow P^{r}$\ldots.}

\uline{Corollaire 1}\qquad $H^{1}(X,\underline{\mathcal O}_{X})\longrightarrow H^{1}(Y_{\overline K},\underline{\mathcal O}_{Y_{\overline K}})$ est injectif

\deleted{\uline{Corollaire 2}\qquad $\mathrm{Pic}(X)\longrightarrow\mathrm{Pic}(Y_{\overline K})$ est injectif}

\uline{Corollaire 2}\qquad $\pi_{1}(Y_{\overline K})\longrightarrow\pi_{1}(X)$ est surjectif [c'est \uuline{Bertini}]

\uline{Corollaire \deleted{\unsure{2|3}}3}\qquad $H^{1}(X,\alpha_{p})\unsure{-|\to}H^{1}(X_{\overline K},\alpha_{p})$ injectif

\uline{Corollaire \deleted{\unsure{3|4}}4}\qquad $H^{1}(X,\mu_{p})\unsure{-|\to}H^{1}(Y_{\overline K},\mu_{p})$ injectif.

On part du corollaire 2, connu grâce à Bertini, \uline{ou} Lefschetz-Grauert [utiliser \deleted{\unsure{disque|doit}} $\operatorname{prof}X\geq2$]. Les en déduit $H^{1}(X,\underline{\mathcal O}_{X})^{F}\Rightarrow H^{1}(X\deleted{\unsure{,}}\mathbb Z/p\mathbb Z)\longrightarrow H^{1}(Y_{\overline K},\underline{\mathcal O}_{Y_{\overline K}})^{F}$ est injectif, on est donc ramené pour le cor. 1 de prouver que $H^{1}(X,\underline{\mathcal O}_{X})^{1-F}-H^{1}(Y_{\overline K},\underline{\mathcal O}_{Y_{\overline K}})^{1-F}$ est injectif [N.B., en tous les \unsure{;} le noyau est stable sous $F$\ldots] i.e. au corollaire 3. Donc on est ramené \deleted{\unsure{au}}\inserted{[en car. $p>0$]} \deleted{\unsure{à}} au cas des cor. 3. et 4. En car $p=0$,

% archive sheet number: 98

\scanpage{100}
\deleted{\unsure{provient|provenant} de $H^1(U,G)$, pourvu que $U$ soit \inserted{pris} assez petit.}

On utilisera un argument de spécialisation.

Or en \deleted{\unsure{vertu de}} \inserted{suivant} Cartier, des deux suites exactes
\[
\begin{aligned}
0\to\mathcal O_X\xrightarrow{F}\mathcal O_X\to\mathcal O_X/\mathcal O_X^p\to0\\
1\to\mathcal O_X^*\xrightarrow{F}\mathcal O_X^*\to\mathcal O_X^*/\mathcal O_X^{*p}\to1
\end{aligned}
\]
\margin{($X$ quelconque !,}

\deleted{\unsure{on ti|on tire}} on tire
\[
\begin{aligned}
H^1(X,\alpha_p)&=H^1(X,\mathcal O_X)^{1-F}\simeq H^0(X,\mathcal O_X/\mathcal O_X^p)\xrightarrow{df}H^0(X,\widetilde{\underline{\Omega}}^1_X)\\
\unsure{{}_pH^1|pH^1|H^1}(X,\mu_p)&=H^1(X,\mathcal O_X)^{1-F}\simeq H^0(X,\mathcal O_X/\mathcal O_X^*)\xrightarrow{\frac{df}{f}}H^0(X,\widetilde{\underline{\Omega}}^1_X)
\end{aligned}
\]
et si $X$ est \uline{normal} les deux flèches écrites sont injectives [\deleted{\unsure{leurs|tous}} \inserted{en tous cas,} les noyaux s'identifient, soit aux \uline{diviseurs additifs} sur $X$ annulés par $\unsure{\sqrt{p}|F}$, soit aux diviseurs multiplicatifs annulés par $F$].

Donc on est ramené au

\uline{Lemme} \ $H^0(X,\widetilde\Omega^1_{X/k})\longrightarrow H^0(Y_{\overline K},\widetilde\Omega^1_{Y_{\overline K}/\overline K})$ est injectif

La vérification est presque triviale \ldots

\uline{Corollaire} \ $\underline{\mathrm{Pic}}^\tau_{X_{\overline K}/\overline K}\longrightarrow\underline{\mathrm{Pic}}^\tau_{\deleted{\unsure{X}}Y_{\overline K}/\overline K}$ a un noyau qui est un $p$-\uline{groupe fini unipotent},
% slanted
\margin{purement infinitésimal}
[le fait qu'il soit unipotent résulte déjà de $\operatorname{prof.}X\geq2$, grâce au \deleted{Bertini} Lefschetz-Grauert-Bertini]. ``sans partie fixe''. Si $X$ est normal (pas néc. de $\dim\geq2$), \inserted{$\underline{\mathrm{Pic}}_{X/k}$ est propre sur $k$.} Le ``sans partie fixe'' provient de la normalité; elle implique que l'on a une \uline{immersion fermée} $\underline{\mathrm{Pic}}^\tau_{X/k}\longrightarrow\prod_i\underline{\mathrm{Pic}}_{Y_i/k}$, les $Y_i$ des sections hyperplanes égalt normales. Donc pour prouver
% archive sheet number: 99

\scanpage{101}
que $\underline{\mathrm{Pic}}^\tau_{X/k}$ est propre, on est ramené à le prouver pour une \uline{courbe} normale, où c'est connu. Alors un sous-grp. unipotent est néc. fini, ce qui achève de prouver l'assertion \ldots

\medskip\centerline{\rule{1.2cm}{0.4pt}}

\uline{Démonstration directe et indépendante de la \unsure{car.|con.} du fait que si $X$ est} \inserted{\uline{géom$^t$}} \uline{normal et projectif} $/k$, \deleted{\unsure{le grp.}}\ $\underline{\mathrm{Pic}}^0_{X/k}$ est propre.

On utilise le théorème de structure des groupes algébriques commutatifs, on est ramené à prouver que \unsure{tout|toute} morphisme $\operatorname{Spec}(k[t,t^{-1}])\to\underline{\mathrm{Pic}}^0_{X/k}$ est constant, \deleted{et il suffit [$k$ clos alg. bien \inserted{sûr.}] de prouver que l'application \unsure{sous-jacente|sans augment} est trivial. \unsure{Cela|Ceci} provient de ceci}

\uline{Lemme} \ Supposons $X$ préschéma loc noeth - normal \inserted{existe un fermé $Y$ de codim $\geq2$ \deleted{\unsure{dont|de}} tel que $U:X-Y$ soit régulier}
\margin{Certainement inutile, en procédant de façon convenable}
Alors \deleted{pour tout}
\[
\mathrm{Pic}(X)\longrightarrow\mathrm{Pic}\,X[t,t^{-1}]
\]
est \deleted{bijectif} \uline{surjectif}.

\deleted{C'est injectif, car $X[t]$ a une section sur $X$ (section nulle).} Soit donc $\underline L$ faisceau inversible sur $X'=X[t,t^{-1}]$, \deleted{soit $U$ l'ens. des pts réguliers de $X$} \inserted{on \unsure{sait} que} \inserted{posant $U'=X'(U)$} $\underline L\mid U'\deleted{\unsure{[t]}}$ provient d'un faisceau inversible $\underline L_0$ sur $U$, qui est défini à l'aide d'un diviseur $D_0$ sur $U$.

Son adhérence $D_1=\overline D_0$ est tel que son image inverse sur $X'\deleted{\unsure{[t]}}$ \deleted{\unsure{est}} soit loc. principale. Pour raison de platitude, il en est de $=$ de $D_1$, d'où facilement la conclusion.
% archive sheet number: 100

\scanpage{102}
[ N.B. \quad \uline{Lemme} \ $X$ préschéma, $U$ ouvert rétrocompact, $f:X'\to X$ morphisme \uline{fid plat} et \uline{qu cpt}, $U'=f^{-1}(X')$, On suppose $\boxed{i'_*(\underline{\mathcal O}_{U'})\xleftarrow{\sim}\underline{\mathcal O}_{X'}}$. (ce qui est \uline{équivalent} à la formule $\boxed{i_*(\deleted{\unsure{U}}\underline{\mathcal O}_U)\xleftarrow{\sim}\underline{\mathcal O}_X}$).

\diagram{\begin{tikzcd}[ampersand replacement=\&,column sep=small,row sep=small]
\underline E' \& {}\\
X'\arrow[d,"f"] \& U'\arrow[l,"i'"']\arrow[d,"g"]\\
X \& U\arrow[l,"i"']\\
{} \& \underline F\unsure{.|{}_{0}}
\end{tikzcd}}

Soit $P_{\unsure{m|n}}(X)$ l'ens des classes, à isom près, de faisceaux loc. libres \deleted{\unsure{des|de}} type fini sur $X$, définissons de $\overline{\mathrm m}$ $P(X'),P(U),P(U')$. Alors le diagramme

\diagram{\begin{tikzcd}[ampersand replacement=\&]
P(X)\arrow[r]\arrow[d] \& P(X')\arrow[d]\\
P(U)\arrow[r] \& P(U')
\end{tikzcd}}
% slanted
\margin{N.B. Inutile de se borner à des faisceaux loc libres. On peut \unsure{énoncer|écrire|encore|prendre}, en prenant les faisceaux $\underline E$ tels que $\operatorname{prof} E_x\geq2$ pour $x\in X-U$, et $\underline E'$ en conséquence, \ldots\ [en loc. noeth]}

\uline{cartésien} \deleted{\unsure{cartésien|cocartésien|cohérent}}. En particulier, le diagramme
\diagram{\begin{tikzcd}[ampersand replacement=\&]
\mathrm{Pic}(X)\arrow[r]\arrow[d] \& \mathrm{Pic}(X')\arrow[d]\\
\mathrm{Pic}(U)\arrow[r] \& \mathrm{Pic}(U')
\end{tikzcd}}
est \uline{cartésien}.

En effet, \deleted{\unsure{si}} le foncteur $\underline E\mapsto i^*(\underline E)$ des $\underline E$ loc libres sur $X$ dans les Mod loc libres sur $U$ est \uline{pleinement fidèle}, donc $\deleted{\unsure{P}}P(X)\to P(U)$ est injectif, [et de $\overline{\mathrm m}$ $P(X')\to P(U')$ l'est] donc $P(X)\to P(U)\times_{P(U')}P(X')$ est injectif. \deleted{Soient d'\unsure{autre|abord}} \inserted{Prouvons la} surjectivité. Soient $\underline F$ loc. libre sur $U$, $\underline E'$ loc libre sur $X'$, tels que $i'{}^*(\underline E')\simeq g^*(\underline F)$. Je dis que $\underline E=i_*(\underline F)$ est loc libre $\underset{\substack{f\ \text{fid plat}\\\text{qu. cpt}}}{\Leftarrow} f^* i_*(\underline F)$ est loc. libre, Or
\[
f^*i_*(\underline F)\underset{f\ \mathrm{plat}}{\simeq} i'_*g^*(\underline F)\underset{g^*(\underline F)\simeq i'{}^*(\underline E')}{\simeq} i'_*i'{}^*(\underline E')\underset{i'_*(\underline{\mathcal O}_{U'})\simeq\underline{\mathcal O}_{X'}}{\simeq}\underline E',\quad\text{OK.}
\]
On a $i^*(\underline E)\simeq\underline F$, Je dis que
\[
 f^*(\underline E)\simeq\underline E'\quad \Longleftarrow\quad
 \underbrace{i'{}^*f^*(\underline E)\simeq i'{}^*(\underline E')}\longrightarrow
 g^*i^*(\underline E)\simeq g^*(\underline F)
 \quad\text{O.K.}
\]
% archive sheet number: 101

\scanpage{103}
\uline{Corollaire} \ Si $\mathrm{Pic}(\deleted{\unsure{X}}U)\unsure{\text{---}|\longrightarrow}\mathrm{Pic}(U')$ est injectif (\unsure{surj|ou|resp.}, bijectif) Il en est de $\overline{\mathrm m}$ de $\mathrm{Pic}(X)\longrightarrow\mathrm{Pic}(X')$ \ldots\ .
% archive sheet number: 102

\scanpage{104}
Th. de Seshadri et p$^{\mathrm{bs}}$ correspondants
% archive sheet number: 103

